\documentclass[11pt]{article}
\usepackage[T1]{fontenc}
\usepackage[utf8]{inputenc}
\usepackage{lmodern,amsmath,amssymb,amsthm,microtype}
\usepackage[margin=25mm]{geometry}
\usepackage{hyperref}
\hypersetup{hidelinks,pdftitle={A finite-energy mixed operator over arbitrary complete fields}}
\setlength{\emergencystretch}{3em}
\newtheorem{theorem}{Theorem}
\newtheorem{lemma}[theorem]{Lemma}
\newtheorem{proposition}[theorem]{Proposition}
\newtheorem{corollary}[theorem]{Corollary}
\newcommand{\PSH}{\operatorname{PSH}}
\newcommand{\MA}{\operatorname{MA}}
\newcommand{\ddc}{dd^c}
\newcommand{\an}{\mathrm{an}}
\newcommand{\E}{\mathcal E}
\newcommand{\CPSH}{\operatorname{CPSH}}
\title{A finite-energy mixed operator\\over arbitrary complete fields}
\author{P16.4 research continuation 217\\Unreviewed partial result}
\date{7 October 2026}
\begin{document}
\maketitle
\begin{abstract}
We propose a direct extension of the genuine continuous mixed
Monge--Ampere operator to actual finite-energy psh metrics on a smooth
projective variety over any complete characteristic-zero
non-Archimedean field. The dimension is arbitrary and trivial,
discrete and nondiscrete valuations are included. Ordered continuous
energy gaps control dense scalar tests; polarization then gives a
positive mixed limit. An averaging estimate proves integrability of
every finite-energy test. A second averaging estimate controls the
error of unbounded tests uniformly along decreasing nets, yielding
joint weighted continuity and the genuine extended-energy formula.
All potentials are actual all-point psh endpoints. The paper supplies
no general-field equation existence theorem, strict locality or
comparison, and assumes no envelope continuity. Every new proof
awaits independent verification.
\end{abstract}

\section{The actual cone and the proposed operator theorem}

Let $K$ be a complete characteristic-zero non-Archimedean real-valued
field, $X/K$ smooth, geometrically integral and projective of dimension
$n\geq1$, and $L$ an actual ample rational line bundle. Write
$T=X^{\an}$ and $V=(L^n)>0$. The space $T$ is compact Hausdorff.
Fix a continuous psh metric $\psi$ on $L$. The term psh means exactly
the decreasing-Fubini--Study-net definition of
\cite[Definition 7.2]{BE} on the fixed rational bundle. We exclude the
identically $-\infty$ metric. Scalar potentials $u$ mean absolute
metrics $\psi+u$ evaluated at every point, including supported points;
$u=-\infty$ is a literal value. Such potentials are upper
semicontinuous and bounded above. Denote the continuous ones by
$\CPSH$ and the actual singular cone by $\PSH$.

For continuous potentials put $\theta_u=\ddc(\psi+u)$ and use the
genuine continuous mixed measures of \cite[Theorem 8.4]{BE}. Wedges
below are measures on $T$ with the fundamental class of $X$ implicit.
Define
\begin{equation}\label{eq:energy}
 \begin{split}
 E(u)&=\frac1{n+1}\sum_{j=0}^n
             \int_T u\,\theta_u^j\theta_0^{n-j}
                     \quad(u\in\CPSH),\\
 E(u)&=\inf\{E(v):v\in\CPSH,\ v\geq u\}
                     \quad(u\in\PSH),\\
 \E^1&=\{u\in\PSH:E(u)>-\infty\}.
 \end{split}
\end{equation}
Continuous majorants exist by the defining approximation. The
extension takes values in $[-\infty,\infty)$ and retains the actual
potential as its argument. It does not define any measure by naming
a prescribed target.

\begin{theorem}\label{thm:main}
For $u_1,\ldots,u_n\in\E^1$ there is a canonical positive Radon
measure $\MA(u_1,\ldots,u_n)$ of raw mass $V$. It agrees with the
genuine continuous mixed operator, is symmetric and affine in the
curvature of each slot, and is the weak limit along every simultaneous
continuous psh decreasing approximation of the actual endpoints.
No support restriction is imposed. Write $\MA(u)=\MA(u,\ldots,u)$.

Every $w\in\E^1$ is integrable against every such mixed output. On
any common directed index set, if $u_{r,i}\downarrow u_r$ for
$1\leq r\leq n$ and $w_i\downarrow w$, all in $\E^1$ with the
indicated endpoints in $\E^1$, then
\begin{equation}\label{eq:weighted}
 \begin{split}
 \MA(u_{1,i},\ldots,u_{n,i})&\longrightarrow
                   \MA(u_1,\ldots,u_n)\quad\hbox{weakly},\\
 \int_T w_i\,\MA(u_{1,i},\ldots,u_{n,i})&\longrightarrow
                   \int_T w\,\MA(u_1,\ldots,u_n).
 \end{split}
\end{equation}
The genuine energy identity is
\begin{equation}\label{eq:sing-energy}
 E(u)=\frac1{n+1}\sum_{j=0}^n
        \int_Tu\,\MA(\underbrace{u,\ldots,u}_{j},
                       \underbrace{0,\ldots,0}_{n-j}).
\end{equation}
For $u,v\in\E^1$, the actual segment $u_t=(1-t)u+tv$ belongs to
$\E^1$, its energy is a polynomial of degree at most $n+1$, and
\begin{equation}\label{eq:derivative}
 \frac{d}{dt}E(u_t)=\int_T(v-u)\,\MA(u_t)
                      \qquad(0\leq t\leq1).
\end{equation}
At endpoints the derivative has its one-sided meaning. The difference
in this integral is evaluated on a common finite-value full-measure
set, without subtracting infinite values.

Changing the continuous psh reference leaves the absolute finite-energy
class and operator unchanged. For every complete valued extension
$F/K$, pullback preserves energy and finite energy, and the projection
$p:X_F^{\an}\to T$ satisfies
\begin{equation}\label{eq:basechange}
 p_*\MA(u_{1,F},\ldots,u_{n,F})=\MA(u_1,\ldots,u_n).
\end{equation}
The raw mass remains $V$; no field-degree factor is introduced.
\end{theorem}

The theorem is an operator statement. It asserts neither attainment
for arbitrary finite-energy targets nor uniqueness of solutions. Its
original P16.4 scope audit is given at the end.

\section{Continuous estimates and compact Radon facts}

The precise continuous inputs are the actual psh closure and smooth
regularization in \cite[Section 7]{BE}, the complete continuous
mixed-operator construction in \cite[Section 8.1]{BE}, and continuous
energy in \cite[Definition 9.13 and Proposition 9.14]{BE}.
In particular the relative cocycle gives, for continuous $u,v$,
\begin{equation}\label{eq:energy-gap}
 E(v)-E(u)=\frac1{n+1}\sum_{j=0}^n
                    \int_T(v-u)\,\theta_v^j\theta_u^{n-j}.
\end{equation}
Continuous energy is monotone, satisfies $E(u+c)=E(u)+cV$, and is
Lipschitz in uniform distance with constant $V$.

Here and below positive real convex combinations of continuous fixed
bundle psh metrics are allowed. For rational coefficients the tensor
product of defining FS systems gives the assertion, followed by
uniform closure. Approximating real coefficients by rationals gives
the continuous assertion. The continuous mixed operator extends the
same affine curvature identity by uniform continuity.

For any product $Q$ of $n-1$ continuous psh curvatures and for scalar
differences of continuous psh metrics, define
\[
 D_Q(f,g)=-\int_T f\,\ddc g\,Q.
\]
The differences may lie on other semiample rational bundles. This is
the signed dpsh calculus of \cite[Section 8.1]{BE}. The form is
symmetric, positive semidefinite and satisfies Cauchy--Schwarz.
To check the latter assertions, uniformly regularize the finitely
many metrics by smooth semipositive metrics using Theorem 7.14 of
that source. For smooth data Stokes and positivity give
$D_Q(f,f)=\int df\wedge d^cf\,Q\geq0$ and symmetry. Analytification
is boundaryless and projectivity supplies compact support for this
Stokes application. The continuous mixed signed measures have
uniformly bounded total variation, bounded by the sum of their
positive intersection masses. Thus for $f_a\to f$ uniformly and
signed measures $\lambda_a\to\lambda$ weakly,
\[
 \left|\int(f_a-f)\lambda_a\right|
       \leq\|f_a-f\|_\infty\|\lambda_a\|_{\rm TV}\longrightarrow0,
 \qquad \int f(\lambda_a-\lambda)\longrightarrow0.
\]
This passes symmetry and positivity to $D_Q$. Positivity of
$D_Q(f+sg,f+sg)$ for real $s$ gives Cauchy--Schwarz. All expressions
are the representation-independent continuous dpsh operator; no
singular measure is used here.

It also proves continuous energy concavity. Along a continuous psh
segment with difference $d=v-u$, its second derivative is
\[
 \frac{d^2}{dt^2}E((1-t)u+tv)
       =-nD_{\theta_{(1-t)u+tv}^{n-1}}(d,d)\leq0.
\]
The first derivative and the polynomial identity follow from the
continuous energy calculus and multilinearity. Thus concavity is
derived from smooth positivity, rather than attributed to a property
not listed in Proposition 9.14.

\begin{lemma}\label{lem:weak-gap}
For every $f\in\mathrm{PL}_{\mathbb R}(T)$ there is $C_f\geq0$,
independent of continuous fixed-$L$ psh $u,v$, such that $u\leq v$
implies
\begin{equation}\label{eq:weak-gap}
 \left|\int_T f\,(\theta_v^n-\theta_u^n)\right|
       \leq\sqrt{n(n+1)C_f\,[E(v)-E(u)]}.
\end{equation}
These real PL functions are uniformly dense in $C(T)$, also at a
trivial valuation.
\end{lemma}
\begin{proof}
By \cite[Lemma 5.17 and Theorem 5.18]{BE}, write $f=h-k$ for two
continuous FS metrics on a common ample bundle $A$; their theorem
with real weights gives the stated density without requiring a
nontrivial absolute value. Put
$C_f=2\|f\|_\infty(A\cdot L^{n-1})$.
Let $\delta=v-u\geq0$ and
$Q_j=\theta_u^j\theta_v^{n-1-j}$ for $0\leq j<n$.
The two positive measures $\ddc h\,Q_j$ and $\ddc k\,Q_j$ have
mass $(A\cdot L^{n-1})$, so
\[
 0\leq D_{Q_j}(f,f)\leq C_f.
\]
Factoring the curvature difference and integrating by parts gives
\[
 \begin{split}
 \sum_{j=0}^{n-1}D_{Q_j}(\delta,\delta)
   &=\int_T\delta(\theta_u^n-\theta_v^n)\\
   &\leq\int_T\delta\theta_u^n
    \leq(n+1)[E(v)-E(u)],
 \end{split}
\]
where the final inequality is the $j=0$ term of
\eqref{eq:energy-gap}; every term there is nonnegative.
Consequently the absolute value of
$\sum_jD_{Q_j}(f,\delta)=-\int f(\theta_v^n-\theta_u^n)$
is at most
\[
 \sum_j\sqrt{D_{Q_j}(f,f)D_{Q_j}(\delta,\delta)}
 \leq\sqrt{nC_f\sum_jD_{Q_j}(\delta,\delta)},
\]
which proves \eqref{eq:weak-gap}.
\end{proof}

\begin{lemma}\label{lem:energy-net}
If actual psh $u_i\downarrow u$ and $u$ is actual psh, then
$E(u_i)\downarrow E(u)$, including the value $-\infty$.
Finite positive convex combinations of actual finite-energy
potentials are actual finite-energy potentials, and extended energy
is concave on this cone.
\end{lemma}
\begin{proof}
Monotonicity gives a decreasing limit at least $E(u)$. For any
continuous majorant $v\geq u$ and $\epsilon>0$, the closed compact
sets $\{u_i\geq v+\epsilon\}$ decrease with empty intersection.
The finite intersection property gives an index after which
$u_i<v+\epsilon$ everywhere. Hence $E(u_i)\leq E(v)+\epsilon V$
on this tail. Infimizing over $v$ and then $\epsilon$ gives the reverse
inequality, including $E(u)=-\infty$.

The same compact argument applies when a decreasing net has
identically $-\infty$ limit: each constant potential $-M$ is then
a continuous majorant of that limit, and the tail lies below
$-M+\epsilon$. Its energies consequently tend to $-\infty$ as
$M\to\infty$.

For finitely many actual endpoints choose their continuous decreasing
psh nets. Positive convex combinations form a product-directed
continuous decreasing net with the asserted all-point limit.
The actual psh decreasing-closure property in
\cite[Lemma 7.4]{BE} puts that limit in the possibly identically
$-\infty$ decreasing-closure cone. Continuous concavity bounds its
approximating energies below by a finite convex combination of the
endpoint energies. The preceding constant-majorant argument excludes
an identically $-\infty$ limit. It is therefore in the actual cone.
Continuous concavity and the just-proved energy limit give
$E(\sum_r a_ru_r)\geq\sum_r a_rE(u_r)>-\infty$ for
$a_r\geq0$, $\sum a_r=1$. The same passage proves concavity.
Zero coefficients are omitted, so $0\cdot(-\infty)$ is never used.
\end{proof}

\begin{lemma}\label{lem:radon}
Let $\mu$ be a finite Radon measure on a compact Hausdorff space.
For a directed increasing net of nonnegative lower semicontinuous
functions $g_i$ with supremum $g$,
\begin{equation}\label{eq:lsc-net}
 \int g\,d\mu=\sup_i\int g_i\,d\mu.
\end{equation}
If finite positive Radon measures $\mu_a\to\mu$ weakly, then for
every nonnegative lower semicontinuous $g$,
\begin{equation}\label{eq:portmanteau}
 \int g\,d\mu\leq\liminf_a\int g\,d\mu_a.
\end{equation}
\end{lemma}
\begin{proof}
For each $t\geq0$, the open set $\{g>t\}$ is the union of the
directed open sets $\{g_i>t\}$. Inner regularity and finite subcovers
of compact subsets imply
$\mu\{g>t\}=\sup_i\mu\{g_i>t\}$. Approximate the layer-cake
integral of $g$ from below by a finite lower step sum, truncating its
height first. For its finitely many levels choose one common index
making all corresponding open-set masses as close as desired to
their suprema. The same lower step sum for $g_i$ proves
\eqref{eq:lsc-net}, even if the integral is infinite.

For \eqref{eq:portmanteau}, compact Hausdorff normality and Radon
inner regularity give
$\int g\,d\mu=\sup\{\int h\,d\mu: h\in C(T),\ 0\leq h\leq g\}$.
One direct verification chooses inner compact subsets at finitely
many strict open levels of $g$ and continuous functions supported
in these open sets and equal to one on the chosen compacts. For
equally spaced levels of spacing $\eta$, their sum multiplied by
$\eta$ is at most $g$ at each point. Such sums approximate the
truncated layer integral, then the full integral. Each fixed
continuous minorant $h$ passes through weak convergence and is
bounded above by $g$ for every $\mu_a$. Supremizing gives
\eqref{eq:portmanteau}. No metrizability or arbitrary measurable-net
monotone convergence is used.
\end{proof}

\section{A canonical pure measure and positive polarization}

\begin{proposition}\label{prop:pure}
For every actual $u\in\E^1$, any continuous psh net
$u_a\downarrow u$ has a unique weak curvature limit of raw mass
$V$, independent of the approximation. Denote it by $\MA(u)$.
The estimate \eqref{eq:weak-gap} extends to ordered $u\leq v$ in
$\E^1$, and $u_i\downarrow u$ in $\E^1$ implies
$\MA(u_i)\to\MA(u)$ weakly.
\end{proposition}
\begin{proof}
The energies $E(u_a)$ decrease to the finite number $E(u)$ by
Lemma~\ref{lem:energy-net}. On each fixed real PL test,
Lemma~\ref{lem:weak-gap} makes the continuous measures a Cauchy net:
for two tail indices use an index dominating both and the triangle
inequality. Positivity and their common mass $V$ bound these
functionals on $C(T)$ uniformly. PL density therefore gives a limit
on every continuous test. Its positive bounded linear functional
defines a positive Radon measure by Riesz, with mass $V$.

For two approximation nets $u_a$ and $v_b$ of the same actual
endpoint, $m_{a,b}=\max(u_a,v_b)$ is a continuous psh decreasing
product net with the same endpoint. Its energies also tend to $E(u)$.
The ordered estimate compares each original measure to that of
$m_{a,b}$ with an error tending to zero. Hence both limits agree.
When $u$ is continuous the constant approximation proves agreement
with its genuine continuous measure.

For ordered singular $u\leq v$, choose continuous approximations
$u_a\downarrow u$, $v_b\downarrow v$. Then
$\max(u_a,v_b)\downarrow v$. Apply the continuous ordered estimate
between $u_a$ and this maximum and pass to their weak limits and
finite energy limits. This gives \eqref{eq:weak-gap} for $u,v$.
For $u_i\downarrow u$ in $\E^1$, Lemma~\ref{lem:energy-net}
and that singular estimate make all PL-test differences tend to zero.
Uniform mass and PL density give weak convergence on $C(T)$.
\end{proof}

For a nonempty subset $I\subset\{1,\ldots,n\}$ set
$u_I=|I|^{-1}\sum_{r\in I}u_r$, which belongs to $\E^1$ by
Lemma~\ref{lem:energy-net}. Define
\begin{equation}\label{eq:polarization}
 \MA(u_1,\ldots,u_n)=\frac1{n!}
    \sum_{\varnothing\ne I\subset\{1,\ldots,n\}}
       (-1)^{n-|I|}|I|^n\MA(u_I).
\end{equation}
For continuous inputs this is the genuine mixed measure: expand
$\theta_{u_I}=|I|^{-1}\sum_{r\in I}\theta_{u_r}$. Inclusion--exclusion
cancels every monomial missing one of the $n$ labels; the surviving
product has coefficient $n!$. This also proves
$\sum_{I\ne\varnothing}(-1)^{n-|I|}|I|^n=n!$.

For simultaneous continuous decreasing approximations of the actual
$u_r$, each subset-average measure tends weakly to $\MA(u_I)$.
At every stage the expression in \eqref{eq:polarization} is a positive
continuous mixed measure of mass $V$. Thus its limit is positive and
has mass $V$. This proves positivity, approximation independence
and continuous agreement without assuming positivity of a signed
formula by itself. Passing the continuous identities proves symmetry
and affine curvature dependence in every slot. For arbitrary singular
$u_{r,i}\downarrow u_r$ the subset averages decrease to their actual
finite-energy endpoints; Proposition~\ref{prop:pure} and the finite
formula give mixed weak convergence.

\section{Every finite-energy test is integrable}

A common constant change of reference makes any fixed finite family
of scalar potentials and tails of their defining continuous nets
nonpositive. We may therefore establish integrability under that
normalization; bounded constant changes do not affect it. Write
$B(u)=-E(u)\geq0$ for $u\leq0$.

\begin{proposition}\label{prop:integrability}
For nonpositive $w,u_1,\ldots,u_n\in\E^1$,
\begin{equation}\label{eq:integrability-bound}
 \int_T(-w)\,\MA(u_1,\ldots,u_n)
   \leq \frac{(n+1)^{n+1}}{n!}
                      \left(B(w)+\sum_{r=1}^nB(u_r)\right)<\infty.
\end{equation}
In particular every mixed output gives zero mass to the pole set
of every actual finite-energy test.
\end{proposition}
\begin{proof}
First assume all potentials continuous. Put
$z=(w+u_1+\cdots+u_n)/(n+1)$. Concavity gives
$B(z)\leq[B(w)+\sum_rB(u_r)]/(n+1)$.
Every term in the formula for $B(z)$ is nonnegative. Its top term
$\frac1{n+1}\int(-z)\theta_z^n$ contains the term
\[
 \frac{n!}{(n+1)^{n+2}}
        \int_T(-w)\theta_{u_1}\cdots\theta_{u_n}
\]
with the displayed coefficient: one factor comes from $-z$, $n$
from the curvature average, and one from the energy denominator.
Retaining this nonnegative term proves
\eqref{eq:integrability-bound} for continuous data.

Fix a nonpositive continuous test $w_k$ and let continuous curvature
approximants of the actual $u_r$ decrease to them. Mixed weak
convergence and finite energy convergence give the same inequality
with $w_k$ and the singular $u_r$. Now take nonpositive continuous
$w_k\downarrow w$. Its costs satisfy $B(w_k)\leq B(w)$.
The functions $-w_k$ increase, are nonnegative and continuous,
and have supremum $-w$. Lemma~\ref{lem:radon} on the single mixed
limit measure gives \eqref{eq:integrability-bound}. Since $w$
is bounded above, this proves full absolute integrability. Its
$-\infty$ set has zero measure as asserted.
\end{proof}

No uncountable intersection of full-measure sets is asserted. Each
finite collection of tests has its own common finite-value set of
full measure, which is all that the ensuing formulas require.

\section{The uniform error of a singular test}

\begin{lemma}\label{lem:test-error}
Let $u,w\in\E^1$ and let continuous psh $w_k\downarrow w$.
For every fixed $k$,
\begin{equation}\label{eq:test-error}
 \int_T(w_k-w)\,\MA(u)
 \leq A_n\left[
 E\left(\frac{u+w_k}{2}\right)
       -E\left(\frac{u+w}{2}\right)\right],
 \qquad A_n=2^{n+1}(n+1).
\end{equation}
Both sides are defined and nonnegative. Differences of potentials
inside integrals mean their finite-value representatives.
\end{lemma}
\begin{proof}
First suppose $u$ is continuous and take a tail index $j\geq k$.
Put $h_k=(u+w_k)/2$, $h_j=(u+w_j)/2$; these are continuous psh
and $h_k\geq h_j$. The top term of \eqref{eq:energy-gap} gives
\[
 (n+1)[E(h_k)-E(h_j)]
       \geq\int_T(h_k-h_j)\MA(h_k).
\]
The positive curvature expansion of
$\MA(h_k)=((\theta_u+\theta_{w_k})/2)^n$
gives $\MA(h_k)\geq2^{-n}\MA(u)$. Since
$h_k-h_j=(w_k-w_j)/2\geq0$, this proves
\eqref{eq:test-error} with $w_j$ in place of $w$.
Let $j$ increase along its directed tail. The functions $w_k-w_j$
are increasing nonnegative continuous functions with supremum
$w_k-w$, so Lemma~\ref{lem:radon} passes the left side to its
integral. Lemma~\ref{lem:energy-net} passes $E(h_j)$ to
$E((u+w)/2)$, finite by convexity. This proves the claim for
continuous $u$.

For actual singular $u$, take continuous $u_a\downarrow u$.
The just-proved inequality applies to every $u_a$. The nonnegative
function $g=w_k-w$ is lower semicontinuous, possibly infinite on
the pole set of $w$. It is integrable against all measures involved
by Proposition~\ref{prop:integrability}. Weak convergence
$\MA(u_a)\to\MA(u)$ and \eqref{eq:portmanteau} give a lower bound
for the liminf of these left sides. Both energies on the right
decrease to their finite actual endpoint energies by
Lemma~\ref{lem:energy-net}, so their difference has the asserted
limit. This proves \eqref{eq:test-error} in general.
\end{proof}

\begin{proposition}\label{prop:weighted-pure}
If $u_i\downarrow u$ in $\E^1$ with $u\in\E^1$, then for
every fixed $w\in\E^1$,
\begin{equation}\label{eq:fixed-weight}
 \int_Tw\,\MA(u_i)\longrightarrow\int_Tw\,\MA(u).
\end{equation}
If also $w_i\downarrow w$ in $\E^1$ with $w\in\E^1$ on the
same index set, the same conclusion holds with $w_i$ on the left.
\end{proposition}
\begin{proof}
Choose continuous psh $w_k\downarrow w$. The product-directed
potentials $(u_i+w_k)/2$ decrease to $(u+w)/2\in\E^1$.
For any $\epsilon>0$, Lemma~\ref{lem:energy-net} gives $i_0,k_0$
such that
\[
 0\leq E((u_i+w_k)/2)-E((u+w)/2)<\epsilon
               \quad(i\geq i_0,\ k\geq k_0).
\]
Also $E((u_i+w)/2)\geq E((u+w)/2)$. Hence
Lemma~\ref{lem:test-error} bounds
$\int(w_{k_0}-w)\MA(u_i)$ by $A_n\epsilon$ on the tail,
and gives the same bound for the limit measure. For the fixed
continuous test $w_{k_0}$, weak convergence gives convergence of
its integrals. The triangle inequality and these two error bounds
prove \eqref{eq:fixed-weight} as $\epsilon\downarrow0$.

For the moving weights, fix $k$ and $\eta>0$. The closed compact
sets $\{w_i\geq w_k+\eta\}$ decrease to an empty intersection
because $w\leq w_k$. Thus $w_i\leq w_k+\eta$ on a tail, and
\[
 0\leq\int_T(w_i-w)\MA(u_i)
       \leq\int_T(w_k-w)\MA(u_i)+\eta V.
\]
Choose $k=k_0$ in the preceding uniform bound, then take a common
tail with the compactness bound. The right side is at most
$A_n\epsilon+\eta V$. All differences are legitimate finite
integrals by Proposition~\ref{prop:integrability}. This proves
that the moving-weight error tends to zero, completing the proof.
\end{proof}

Apply this proposition to each subset-average net in
\eqref{eq:polarization}. Its finite signed coefficients are fixed,
and all test integrals exist. Summing therefore proves the full
mixed weighted law \eqref{eq:weighted}. The argument works for
arbitrary directed nets; sequences or a single skeleton were never
substituted for them.

\section{Energy, reference covariance and field extension}

Choose continuous $u_i\downarrow u\in\E^1$. In each term of the
continuous energy identity take $u_i$ as moving weight, $j$ slots
equal to $u_i$ and $n-j$ slots equal to the reference potential zero.
The mixed weighted law and Lemma~\ref{lem:energy-net} give
\eqref{eq:sing-energy}. This proves that the previously defined
majorant energy is the energy of the genuine limiting operators.

For the segment of two actual finite-energy potentials $u,v$,
convexity was proved in Lemma~\ref{lem:energy-net}. Let continuous
$u_a\downarrow u$, $v_b\downarrow v$. At every product-directed
stage the continuous energy on their affine segment is a polynomial
of degree at most $n+1$, and its derivative is the integral of
$v_b-u_a$ against its pure measure. Each coefficient of this derivative
is a fixed binomial factor times an integral with weight $u_a$ or
$v_b$ and curvature slots chosen from $u_a,v_b$. The full mixed
weighted law passes each such coefficient to its finite singular
integral. The constant coefficient of the energy converges to
$E(u)$. There are finitely many coefficients, so integration of
these coefficient limits proves a limiting polynomial whose value
at any $t$ is $E(u_t)$, by decreasing energy continuity. Its derivative
is precisely the polarized integral in \eqref{eq:derivative}.
This is a finite polynomial argument, not an exchange of arbitrary
limits and derivatives. Proposition~\ref{prop:integrability} permits
$v-u$ on the common finite-value full-measure set of each output.

Let $\psi'$ be another continuous psh reference on $L$, and write
$b=\psi'-\psi$ and $u'=u-b$. The set of continuous absolute psh
majorants of $\psi+u=\psi'+u'$ is the same for both references.
The continuous cocycle gives
$E_{\psi'}(v-b)=E_\psi(v)-E_\psi(b)$ for every majorant.
Infimizing gives
\[
 E_{\psi'}(u')=E_\psi(u)-E_\psi(b).
\]
Thus the absolute finite-energy class is identical. Its absolute
continuous approximation nets are also identical; their genuine
curvature measures do not depend on reference. The canonical limits
and all mixed operators are therefore identical. This justifies the
constant normalizations used above as well as general reference
covariance.

For a complete valued extension $F/K$, the standard analytic
projection is continuous and surjective. Pulling back the defining
continuous FS nets gives the actual endpoint $u_F=u\circ p$.
The continuous energies agree under extension by
\cite[Proposition 9.14]{BE}. Decreasing energy continuity on both
spaces gives
\[
 E_{\psi_F}(u_F)=\lim_aE_{\psi_F}(u_{a,F})
              =\lim_aE_\psi(u_a)=E_\psi(u).
\]
In particular $u_F\in\E^1$. The genuine continuous mixed
pushforward law of \cite[Proposition 8.3 and Section 8.1]{BE}
passes to weak limits by testing with $f\circ p$ for $f\in C(T)$,
and proves \eqref{eq:basechange}. Intersection degree is unchanged
under extension. Neither an exact extension-supnorm identity nor
descent of arbitrary upstairs singular metrics was assumed.

\section{Zero dimension, imported scope and remaining equation inputs}

For $n=0$, the retained geometric-integrality hypothesis gives
$X=\operatorname{Spec}K$ and $T$ a single point. The actual nontrivial
psh scalars are finite constants, $V=1$, $E(u)=u$ and the operator
is its unit point mass. All weighted identities and field laws then
hold directly. The substantial argument above concerns $n\geq1$.

The imported continuous theory is used in its precise projective,
geometrically reduced, semiample rational-bundle domain; the present
smooth geometrically integral ample setting meets all those
hypotheses. The real PL density is the source's real-weight theorem,
so it includes trivial valuations. Smooth regularization and the
continuous mixed theory retain their whole carrier and raw
intersection mass. Smooth Stokes/positivity is used with boundaryless
analytification and compact support; the proof supplies no new
recursive reconstruction of the CLD calculus. Compact Radon/Riesz
and normality are ordinary foundations, with the nonmetrizable net
claims proved explicitly above. No accepted earlier partial result
is a mathematical grant for this manuscript. Each new inference
must be independently audited together with the precise primary
continuous-source applications.

The paper constructs a proposed arbitrary-field finite-energy mixed
operator, all-test integrability, weighted decreasing continuity,
energy and affine-segment calculus. It does not yet prove strict
Borel locality or comparison, compactness of normalized actual
potentials, arbitrary-field envelope continuity, orthogonality or
both-sign differentiation in outer continuous directions. It gives
no equivalence between all-test target integrability and a finite
dual cost, no actual maximizing potential for an arbitrary prescribed
target, and no every-solution uniqueness theorem. These remain
necessary for the unchanged original P16.4 equation request. The
original stays open even if this complete partial manuscript passes
independent review. The historical source statement that the general
extension was unclear is not used as either a mathematical premise
or a claim that present literature has been exhaustively searched.

\begin{thebibliography}{9}
\bibitem{BE}
S. Boucksom and D. Eriksson,
\emph{Spaces of norms, determinant of cohomology and Fekete points
in non-Archimedean geometry}, arXiv:1805.01016v3 (2021).
\url{https://arxiv.org/pdf/1805.01016v3}.
\bibitem{CLD}
A. Chambert-Loir and A. Ducros,
\emph{Formes differentielles reelles et courants sur les espaces de
Berkovich}, arXiv:1204.6277v1 (2012).
\url{https://arxiv.org/pdf/1204.6277v1}.
\end{thebibliography}
\end{document}
